\section{The two-degree increment and physical-time transport}
\label{sec:degree}

Write \(G_-=G_N(t,z)\) and \(G_+=G_{N+2}(t,z)\), with \(N\ge2\). For any
quantity evaluated in the two fields, let \(\Delta X=X_+-X_-\).
The two-degree numerator \(d_G=G_+-G_-\) is determined by a finite
coefficient prefix at each time. Its comparison with the heat residual
requires the next link. The finite formula and the argument on a common
Fourier zero contour follow the degree and time identities in \cite{jh03}.

\subsection{A shared finite numerator}

\begin{theorem}[Two-degree resolvent and simplex formula]
\label{d:finite-numerator}
Suppress \(t\) and write \(P_j(\xi)=\det(\xi I-J_j)\) and
\(s_j(\xi)=((\xi I-J_j)^{-1})_{00}\). As in
\eqref{eq:full-phase}, \(L_j=h_j(iz)^{2j}/(2j)!\) with
\(h_j=\prod_{k=1}^j\beta_k\).
Then
\begin{equation}
 s_{N+2}(\xi)-s_N(\xi)=
       \frac{h_N\xi}{P_N(\xi)P_{N+2}(\xi)}.
 \label{d:resolvent}
\end{equation}
Let \(x_1,\ldots,x_m\), \(m=2N+2\), be the combined list of spectral
nodes of \(J_N,J_{N+2}\), with every occurrence retained. Set
\(X=\sum_{i=1}^m u_i x_i\) for \(u\) uniformly distributed on the
probability simplex and let \(\chi(z)=\mathbb E e^{izX}\).
For a positively oriented spectral contour \(\Sigma\) enclosing the nodes,
\begin{align}
 d_G(t,z)
 &=\frac{h_N}{2\pi i}\oint_\Sigma
           \frac{\xi e^{iz\xi}}{P_N(\xi)P_{N+2}(\xi)}\,d\xi
 \label{d:contour-numerator}\\
 &=L_N\left(\chi(z)+\frac{z\chi'(z)}{2N+1}\right).
 \label{d:simplex-numerator}
\end{align}
In particular, the parenthesized factor is entire, including at
\(z=0\). For \(y_*\ge|\operatorname{Im}z|\) and
\(K\ge\|J_{N+2}\|\), its deviation from \(1\) satisfies
\begin{equation}
 \left|\chi+\frac{z\chi'}{2N+1}-1\right|
 \le |z|^2\left(\frac12+\frac1{2N+1}\right)
     \sigma^2 e^{y_*K},
 \qquad
 \sigma^2=
 \frac{\operatorname{tr}J_N^2+\operatorname{tr}J_{N+2}^2}
      {(2N+2)(2N+3)} .
 \label{d:simplex-bound}
\end{equation}
\end{theorem}

\begin{proof}
Set \(C_j(\xi)=\xi I-J_j\) and write \(a_j=\sqrt{\beta_j}\).
Appending the last site changes the \(00\) resolvent entry by
\[
 \frac{\beta_j(C_j^{-1})_{0,j-1}^{\,2}}
      {\xi-\beta_j(C_j^{-1})_{j-1,j-1}}.
\]
The endpoint cofactor gives
\((C_j^{-1})_{0,j-1}=\prod_{k=1}^{j-1}a_k/P_j\).
The determinant formula for the Schur complement gives the denominator
\(P_{j+1}/P_j\). Hence \(s_{j+1}-s_j=h_j/(P_jP_{j+1})\).
Adding consecutive differences and expanding the last row and column of
\(P_{N+2}\) gives \(P_{N+2}+\beta_{N+1}P_N=\xi P_{N+1}\), which
proves \eqref{d:resolvent} as a rational identity. Integrating the
scalar resolvents against \(e^{iz\xi}\), or applying Cauchy's formula
in an eigenbasis, gives \eqref{d:contour-numerator}.

For the simplex factor put \(p=m-1\), and integrate in the \(p\) free
coordinates \(u_i\ge0\), \(\sum_{i=1}^{p}u_i\le1\), with
\(u_m=1-\sum_{i=1}^p u_i\). The volume is \(1/p!\). For a multi-index
\(\alpha\) with \(|\alpha|=k\), successive beta integrals give
\(\mathbb E\prod_i u_i^{\alpha_i}=p!\prod_i\alpha_i!/(p+k)!\).
Consequently the multinomial formula gives
\(\mathbb E X^k=k!p!H_k(x)/(p+k)!\), where
\(H_k(x)=\sum_{|\alpha|=k}\prod_i x_i^{\alpha_i}\).
On a large circle the Laurent expansion
\(\prod_i(\xi-x_i)^{-1}=\xi^{-m}\sum_{k\ge0}H_k(x)\xi^{-k}\)
therefore gives the Hermite--Genocchi formula
\[
 \frac1{2\pi i}\oint_\Sigma
      \frac{f(\xi)}{\prod_i(\xi-x_i)}\,d\xi
 =\frac{1}{p!}\mathbb E f^{(p)}(X).
\]
Indeed both sides give \(H_k(x)\) for \(f(\xi)=\xi^{p+k}\), and zero
for monomials of degree less than \(p\). The formula follows for entire
\(f\) by locally uniform Taylor expansion and contour deformation.
This proof also covers repeated nodes. For \(f(\xi)=e^{iz\xi}\), the
right side is \((iz)^p\chi(z)/p!\). Applying \(i^{-1}\partial_z\)
inserts the numerator \(\xi\) and gives \eqref{d:simplex-numerator}.

The node lists have zero sum, so \(\mathbb E X=0\). The simplex
second moments are \((1+\mathbf 1_{i=j})/[m(m+1)]\), giving the stated
variance. Compression and convexity give \(|X|\le K\). Taylor's
integral remainders, using \(\mathbb E X=0\), give
\[
 |\chi-1|\le \tfrac12 |z|^2\mathbb E[X^2e^{y_*|X|}],
 \qquad
 |\chi'|\le |z|\mathbb E[X^2e^{y_*|X|}],
\]
which imply \eqref{d:simplex-bound}.
\end{proof}

The poles of the resolvents in \(\xi\) are nodes of the finite source
quadratures. We count zeros of the entire Fourier readouts on a
separate contour in \(z\). Formulas \eqref{d:contour-numerator}
and \eqref{d:simplex-bound} use links through \(\beta_{N+1}\).
By \eqref{eq:finite-defect}, the two normalized heat residuals use
through \(\beta_{N+2}\).

The full-source error laws \eqref{eq:full-phase} give a second
expression for \(d_G\). Put
\(A_j=\beta_{j+1}z^2/[(2j+1)(2j+2)]\), so \(L_{j+1}=-A_jL_j\).
Subtraction at the two degrees gives
\begin{equation}
 d_G=L_N(E_N-A_NA_{N+1}E_{N+2}).
 \label{d:source-numerator}
\end{equation}
An enclosure of \(E_j-1\) by a disk of radius \(e_j\) puts the
parenthesized factor in the disk with center \(1-A_NA_{N+1}\) and
radius \(e_N+|A_NA_{N+1}|e_{N+2}\). The explicit coefficients in
\eqref{d:source-numerator} use through \(\beta_{N+2}\). Evaluating
each exact \(E_j\) requires the full source. The inherited phase
enclosures use their own strip and time hypotheses
\cite{vcd,jh03,jh04}.

\subsection{A single logarithm for both counted coordinates}

Fix \(\Gamma:|z-a|=r\) with real \(a\), oriented positively. On a
compact time interval \(I\), assume each field has two zeros inside
\(\Gamma\), with multiplicities, and the fields share a zero-free
annular collar. Assume both pairs lie in \(|z-a|\le r-b\), where
\(0<b\le r\). Write \(B_j=G_{j,z}/G_j\), \(w_j=z-c_j\),
\(Q_j=w_j^2-D_j\), \(j\in\{-,+\}\). On \(\Gamma\) the individual
root distances are at least \(b\). Thus
\[
 |Q_j|\ge b^2,\qquad |w_j|\le g:=2r-b,\qquad
 |D_j|\le (r-b)^2.
\]
Equal counts give zero winding for \(G_+/G_-\) on the collar, so it
has a holomorphic logarithm \(\ell\). A choice at one time continues
in \(I\) by integration of the time logarithmic derivative on the
collar. All the contour identities are independent of the constant
branch choice.

\begin{proposition}[Logarithmic coordinate difference]
\label{d:log-coordinates-prop}
With \(\bar c=(c_++c_-)/2\),
\begin{equation}
 \Delta c=-\tfrac12\langle\ell\rangle,\qquad
 \Delta D=-\langle(z-c_-)\ell\rangle-(\Delta c)^2
         =-\langle(z-\bar c)\ell\rangle.
 \label{d:coordinate-log}
\end{equation}
If \(L\ge\sup_\Gamma|\ell|\) for some branch, then
\(|\Delta c|\le rL/2\) and \(|\Delta D|\le rgL\).
The formulas extend through a double zero inside the contour.
\end{proposition}

\begin{proof}
For \(p_{k,j}=\langle(z-a)^kB_j\rangle\), the argument principle
gives \(c_j=a+p_{1,j}/2\), \(D_j=p_{2,j}/2-p_{1,j}^2/4\).
Since \(\Delta B=\ell_z\), integration by parts gives
\(\Delta p_k=-k\langle(z-a)^{k-1}\ell\rangle\). Substitution for
\(k=1,2\) gives \eqref{d:coordinate-log}. The norm of a contour
integral is bounded by \(r\) times the supremum on \(\Gamma\), and
\(|z-\bar c|\le g\). Both formulas use counted moments that are analytic
in time on the cleared chart.
\end{proof}

For example, if \(\sup_\Gamma|d_G/G_-|\le\eta<1\), the series
branch of \(\log(1+d_G/G_-)\) gives \(L=-\log(1-\eta)\).
If the straight interpolation \(G_-+s d_G\), \(0\le s\le1\), has
positive clearance on the circle, its log integral in \(s\) gives
\(L=\sup_\Gamma |d_G|/\inf_{s,\Gamma}|G_-+s d_G|\).
These estimates require the indicated denominator floors.

The exterior trace \(M_j=S_{1,j}=-\langle Q_j^{-1}B_j\rangle\)
also changes with degree. More generally, for a fixed-order trace
\(X_j=-\langle k_j B_j\rangle\) of the pair-dependent kernel \(k_j\),
\begin{equation}
 \Delta X=\langle k_{+,z}\ell\rangle
           -\langle(k_+-k_-)B_-\rangle .
 \label{d:kernel-difference}
\end{equation}
This follows by subtracting the two trace integrals and integrating
\(\ell_z\) by parts. For \(M\), the change of the kernel is given by
\[
 Q_+^{-1}-Q_-^{-1}
   =\frac{(2z-c_--c_+)\Delta c+\Delta D}{Q_+Q_-}.
\]
In particular, if \(a_c,a_D\) bound \(|\Delta c|,|\Delta D|\),
then
\begin{equation}
 |\Delta M|\le
 r\left(\frac{2L}{b^3}
        +\frac{(2g a_c+a_D)\sup_\Gamma|B_-|}{b^4}\right).
 \label{d:delta-m-bound}
\end{equation}
Here \(|(Q_+^{-1})_z|\le 2/b^3\) by differentiating its two
linear factors. Both denominators in the kernel change have their
own floors. Formula \eqref{d:kernel-difference} retains the correlation
with \eqref{d:coordinate-log} before taking absolute values.

\subsection{Residuals and endpoint transport}

Let \(f_j=R_j/G_j\), \(V_j=G_{j,zz}/G_j\), and
\(r_j=-\langle w_jf_j\rangle\), the gap force denoted by \(\delta\)
in \eqref{h:pair}.
The common scalar \(\beta_1\) cancels from the time logarithmic
derivative. The two-field heat identity on the collar is
\begin{align}
 \ell_t+\ell_{zz}+(B_++B_-)\ell_z
  &=\Delta f
   =\frac{\Delta R-f_-d_G}{G_+},
 \label{d:two-field-heat}\\
 \ell_t&=\Delta f-\Delta V
   =\frac{\Delta R-d_{G,zz}-(f_--V_-)d_G}{G_+}.
 \label{d:log-time}
\end{align}
Indeed \(V_j=B_{j,z}+B_j^2\) and
\(G_{j,t}/G_j=f_j-V_j-\beta_1\). In the numerator of
\eqref{d:log-time},
\(\Delta R-d_{G,zz}=(\partial_t+\beta_1)d_G\) vanishes to
order at least \(2N\) at the origin by \eqref{d:simplex-numerator}.
This cancellation concerns the Taylor order at \(z=0\).
The residual difference also includes the changed center.
\begin{equation}
 \Delta r=-\langle w_+\Delta f\rangle+\Delta c\,\langle f_-\rangle.
 \label{d:residual-difference}
\end{equation}
Given a floor \(m_+\le\inf_\Gamma|G_+|\) with \(m_+>0\),
\eqref{d:two-field-heat} gives
\(\|\Delta f\|_\Gamma\le(\|\Delta R\|_\Gamma+
\|f_-\|_\Gamma\|d_G\|_\Gamma)/m_+\), where \(\|\cdot\|_\Gamma\)
denotes the circle supremum. Thus
\(|\Delta r|\le r(g\|\Delta f\|_\Gamma+a_c\|f_-\|_\Gamma)\).

For a differentiable collar test \(k\), integration by parts in
\eqref{d:two-field-heat} gives
\[
 \langle k\,\Delta f\rangle
 =\frac{d}{dt}\langle k\ell\rangle
 -\big\langle\big[k_t-k_{zz}+(k(B_++B_-))_z\big]\ell\big\rangle .
\]
The interior contribution is part of the time integral of the residual.
For the complete gap difference the two center terms combine to give
\begin{equation}
 (\Delta D)'=-4\Delta(DM)+\Delta r
   =-\langle (z-\bar c)\ell_t\rangle
           +\bar c'\langle\ell\rangle .
 \label{d:raw-gap}
\end{equation}
The first equality is the difference of the pair equations
\eqref{h:pair}. The second is differentiation of
\eqref{d:coordinate-log}. Both are regular when either \(D_j=0\).

\begin{theorem}[Complete relative forcing on a time subinterval]
\label{d:transport}
For \(u\le v\) in \(I\), set
\[
 W_j(t;v)=\exp\left(-4\int_t^v M_j(s)\,ds\right)
 \qquad(j\in\{-,+\}, t\in[u,v]).
\]
The traces \(M_j\) are real by conjugation symmetry, so both weights
are positive. The full forcing satisfies
\begin{equation}
 \int_u^v W_-(t;v)\,[\Delta r-4D_+\Delta M](t)\,dt
   =\Delta D(v)-W_-(u;v)\Delta D(u).
 \label{d:weighted-gap}
\end{equation}
For \(u<v\), division by \(\int_u^v W_-(t;v)\,dt>0\) gives the
weighted mean of the full forcing on that subinterval. If
\(a_D(t)\ge|\Delta D(t)|\), the absolute value of the integral in
\eqref{d:weighted-gap} is at most \(a_D(v)+W_-(u;v)a_D(u)\).
\end{theorem}

\begin{proof}
Expand \(\Delta(DM)=M_-\Delta D+D_+\Delta M\), so
\((\Delta D)'+4M_-\Delta D=\Delta r-4D_+\Delta M\).
Since \(\partial_t W_-=4M_-W_-\), the integrand in
\eqref{d:weighted-gap} equals \(\partial_t(W_-\Delta D)\).
Evaluation at the endpoints proves the equality, and positivity
of \(W_-\) gives the stated bound.
\end{proof}

The corresponding formulas with each field's own weight are
\[
 \int_u^v W_j(t;v)(2+r_j(t))\,dt=D_j(v)-W_j(u;v)D_j(u).
\]
Subtracting them yields
\(\Delta D(v)-W_+(u;v)\Delta D(u)
 -(W_+-W_-)(u;v)D_-(u)\).
For \(u<v\), each weight at \(u\) is the exponential of \(-4(v-u)\)
times the respective mean of \(M_j\) on \([u,v]\).
For \(a_M\ge|\Delta M|\), the weight difference satisfies
\[
 |W_+(t;v)-W_-(t;v)|
 \le W_-(t;v)\left(\exp\left(4\int_t^v a_M(s)\,ds\right)-1\right).
\]
There is also a complete endpoint identity for the \(M\) row.
Write \(\mathcal F_j=12S_{2,j}+16D_jS_{3,j}\) and denote its
residual contribution in the equation for \(M_j\) by \(\varepsilon_{M,j}\).
The row \(M_j'=2M_j^2-\mathcal F_j+\varepsilon_{M,j}\) in
\eqref{h:low} gives, for
\(A(t)=\exp(-2\int_u^t(M_++M_-))\),
\[
 \int_u^v A(t)(\Delta\varepsilon_M-\Delta\mathcal F)(t)\,dt
   = A(v)\Delta M(v)-\Delta M(u).
\]
Here \(\Delta\mathcal F=12\Delta S_2+
16(D_+\Delta S_3+S_{3,-}\Delta D)\) retains the full change
in the higher traces and in the moving gap.

The absolute bound in Theorem~\ref{d:transport} applies after
integration of the signed sum. Numerical use still requires count
and clearance on the selected chart, bounds for the endpoint log
functionals and enclosure of the weight integral. For example, an
error \(e_P\) in an evaluated integral of \(M_-\) contributes at
most \(\widetilde W_-(e^{4e_P}-1)\) to the error of the evaluated
weight. A direct contour evaluation furnished with the same
logarithm, endpoint data and weight can use the same cancellation.
For comparison with the complete theta field, the logarithm and time
identities still apply on its common counted chart, with \(R_F=0\).
That comparison requires source-error and clearance estimates for
the complete field. A sign conclusion requires an additional signed
estimate on the selected chart.
